\section{An explicit pencil of directions}\label{sec:pencil}

The first ingredient is a short list of vectors on which every quadratic
evaluation that is nonzero somewhere on the list has large support.
Its length must be linear in the
ambient dimension.  We construct the list by evaluating functions in an
additive tower, and give the local calculations and finite algorithm needed
for its uniformity.

\begin{lemma}[Quadratic evaluation pencil]\label{lem:pencil}
For every $0<\lambda_0<1$ and every sufficiently large power of two $r$,
put $q=r^2$.  There is a deterministic algorithm which, given an integer
$d\geq 1$, constructs a list
\[
 u_1,\ldots,u_m\in\F_q^d
\]
with the following properties:
\begin{enumerate}[label=(\roman*)]
 \item $m=O(d)$, and the list contains every standard basis vector of
 $\F_q^d$;
 \item if a homogeneous quadratic polynomial $Q$ in $d$ variables does
 not vanish at every point of the list, then
 \[
   \#\{j:Q(u_j)=0\}\leq\lambda_0m.
 \]
\end{enumerate}
The running time is $2^{O(d)}$ bit operations.  The constants in both bounds
may depend on $r$, which is fixed independently of $d$.
\end{lemma}

In Section~\ref{sec:bit-tests}, these vectors define rank-one directions
for a random walk on a space of matrices. Its spectral estimate needs
only quadratic evaluations that are nonzero somewhere on the list,
exactly the hypothesis in~\textup{(ii)}. The application takes
$d=O(\log(\text{input size}+1))$, so the construction time is polynomial
in the formula size.

The proof first describes the local branches of the tower and their
integral lattices. Their determinant orders bound the conditions imposed
on global functions with controlled poles. We then evaluate a sufficiently
large space of these functions to obtain the quadratic zero bound, and
finally give a finite algorithm with explicit precision and running-time
bounds.

Fix a power of two $r$ sufficiently large that, with $D=r+1$,
\begin{equation}\label{eq:pencil-parameters}
 \frac{2D}{r^2-r}\leq\lambda_0.
\end{equation}
All finite fields used below have characteristic two.  Write $F=\F_q$ and
let $k$ be an algebraic closure of $F$.  Define
\begin{equation}\label{eq:pencil-tower}
 S(Y)=Y^r+Y,\qquad
 f(Y)=\frac{Y^{r+1}}{Y^r+Y}
     =\frac{Y^r}{Y^{r-1}+1}
     =\frac{1}{S(1/Y)},\qquad
 S(x_{i+1})=f(x_i),
\end{equation}
where $x_1$ is transcendental. This is the Garcia--Stichtenoth tower
\cite{GS1996}; the recurrence also appears in
\cite[Example~3.3]{Stichtenoth2001}.  The identities involving inverses are
rational-function identities.  In particular, the middle expression defines
$f$ at zero.  We will prove that the tower through $x_l$ is a field of degree
$N=r^{l-1}$ over $k(x_1)$, and that the monomials
\begin{equation}\label{eq:pencil-power-basis}
 e_{\boldsymbol a}=x_2^{a_2}\cdots x_l^{a_l},
 \qquad 0\leq a_i<r,
\end{equation}
form a basis.  For $l=1$ this is the one-element basis $1$.
The same equations and basis will then define the tower over $F(x_1)$.

\subsection{The two local extensions}

For a formal Laurent series, its order is the least exponent with nonzero
coefficient; the order of zero is $+\infty$.  An element of $k((z))$ is
called integral if it belongs to $k[[z]]$.

Suppose first that the right side of an equation $S(X)=b(z)$ becomes
integral after subtracting $S(g)$, for a specified $g\in k((z))$.  Put
$b_0=b-S(g)\in k[[z]]$.  There are exactly $r$ choices of a constant
$\beta$ satisfying $\beta^r+\beta=b_0(0)$: the derivative of this polynomial
is one, and its roots differ by the elements of $\F_r$.  For each choice,
there is a unique positive-order series $u$ with
$S(u)=b_0-S(\beta)$.  Indeed, if $b_0-S(\beta)=\sum_{n>0}b_nz^n$, then
\begin{equation}\label{eq:pencil-additive-recursion}
 u_n=b_n+\ind_{\{r\mid n\}}u_{n/r}^{r}\qquad(n>0)
\end{equation}
determines its coefficients successively.  Thus the equation splits into
$r$ distinct roots $g+\beta+u$ in $k((z))$.  Its local algebra is a product
of $r$ copies of $k((z))$, with integral lattice $k[[z]]^r$.

The second case is a simple pole.  It is useful to establish the parameter
and degree assertions without first assuming irreducibility.

\begin{lemma}[A simple-pole extension]\label{lem:pencil-simple-pole}
Let $b\in k((z))$ have order $-1$.  The polynomial $X^r+X-b$ is
irreducible over $k((z))$.  Its extension is a Laurent series field
$k((w))$, where $w=1/X$ and $z$ has order $r$ in $w$.  Its integral lattice
has the basis
\[
 1,w,\ldots,w^{r-1}\quad\text{over }k[[z]],
\]
and the discriminant of this basis has $z$-order $2(r-1)$.
\end{lemma}

\begin{proof}
The series $q_0(z)=1/b(z)$ has order one.  Its compositional inverse
$\psi$ exists: the coefficient of each new term is determined by division
by the nonzero linear coefficient of $q_0$.  In a new series field
$k((w))$, set
\begin{equation}\label{eq:pencil-reparameterization}
 z=\psi\left(\frac{w^r}{1+w^{r-1}}\right).
\end{equation}
Substitution gives an injective homomorphism from $k((z))$ to $k((w))$,
and $z=cw^r+\cdots$ for some $c\ne0$.  Moreover $X=1/w$ satisfies
$X^r+X=b(z)$.

The monomials $w^jz^n$, with $0\leq j<r$ and $n\geq0$, have distinct
leading orders $j+rn$, which exhaust the nonnegative integers.  Given a
power series in $w$, subtract its lowest remaining coefficient using the
unique such monomial of that order, and repeat.  The subtractions converge
formally because their orders tend to infinity.  Collecting terms with the
same $j$ gives
\[
 k[[w]]=\bigoplus_{j=0}^{r-1}w^j k[[z]].
\]
The sum is direct: the leading orders of nonzero summands belong to
distinct residue classes modulo $r$.  After multiplying a Laurent series
by a sufficiently large power of $z$, the same argument applies.  Hence
$[k((w)):k((z))]=r$ and $k((w))=k((z))(w)=k((z))(X)$.
This proves irreducibility, and separability follows from the derivative
$S'(X)=1$.

The displayed free module is closed under multiplication, so each of its
elements is integral over $k[[z]]$ by the characteristic polynomial of its
multiplication map.  Conversely, an element of negative $w$-order cannot
satisfy a monic polynomial over $k[[z]]$: its highest power would be the
unique term of smallest order.  Thus $k[[w]]$ is precisely the integral
closure.

The $r$ embeddings send $X$ to $X+c$, $c\in\F_r$, and therefore send $w$
to $w_c=1/(X+c)$.  For distinct $c,c'$, the difference
\[
 w_c-w_{c'}=\frac{c'-c}{(X+c)(X+c')}
\]
has $w$-order two, or extended $z$-order $2/r$.  The evaluation matrix of
$1,w,\ldots,w^{r-1}$ is Vandermonde.  Its determinant squared is the
trace-pairing determinant, so the latter has order
$2\binom r2(2/r)=2(r-1)$.
\end{proof}

Here and below the trace of an element of a finite algebra is the trace
of its multiplication map, and its norm is the determinant of that map.
For a separable field or product of separable fields, passing to a field
containing all embeddings diagonalizes multiplication.  Thus trace is
the sum of the embedded values, and the trace matrix of a basis is the
transpose of its evaluation matrix times that matrix.  This also proves
the discriminant identity used in Lemma~\ref{lem:pencil-simple-pole}.

\subsection{All branches of the tower}

At a pole of $Y$, the difference $f(Y)-Y$ is integral: from
$f(Y)=Y/(1+Y^{1-r})$, its order is at least zero.  Start at infinity with
the parameter $t=1/x_1$.  The first right side has a simple pole, and
Lemma~\ref{lem:pencil-simple-pole} produces a single branch in which
$x_2$ has a simple pole with parameter $1/x_2$.  Repeating the lemma gives
a single branch at every level.  Inductively, the polynomial defining
$x_{i+1}$ is irreducible over this completion of the preceding field,
and hence over that field itself.  This proves the degree and basis
assertions following Equation~\ref{eq:pencil-power-basis}.  It also proves
them over $F(x_1)$, since an equation irreducible after extending constants
to $k$ was already irreducible over $F$.

At a finite point $x_1=a$, put $t=x_1-a$.  If
$a\in\F_r\setminus\{0\}$, then $S(x_1)$ has a simple zero and
$x_1^{r+1}$ is a unit.  Thus the first right side has a simple pole, and
all subsequent steps are again simple-pole extensions.  If
$a\notin\F_r$, the first right side is integral and the extension splits.
An integral residue $\beta\notin\F_r$ has $f(\beta)\ne0$, so no root of
$S(Y)=f(\beta)$ belongs to $\F_r$.  Induction therefore gives only
unshifted integral splittings at such a point.

It remains to describe every branch above zero.  As long as the residues
of $x_1,x_2,\ldots$ are zero, all extensions split without a shift.  Since
$f(Y)$ has order $r$ at zero and $S(Y)$ has order one there,
\[
 \operatorname{ord}_t(x_j)=r^{j-1}
\]
along the branch with zero residues.  Suppose the first nonzero residue
occurs at $x_i$, where $i\geq2$.  It is an element
$\alpha\in\F_r^*$, and
\[
 \operatorname{ord}_t(x_i-\alpha)
 =\operatorname{ord}_t S(x_i)=r^{i-1}.
\]
The numerator $x_i^{r+1}-\alpha^2$ is divisible by $x_i-\alpha$.
Consequently, modulo integral series,
\begin{equation}\label{eq:pencil-first-shift}
 f(x_i)\equiv\frac{\alpha^2}{S(x_i)}
 =S\left(\frac{\alpha^2}{x_{i-1}}\right).
\end{equation}
The equality uses $S(x_i)=f(x_{i-1})$, the last identity in
Equation~\ref{eq:pencil-tower}, and $\alpha^2\in\F_r$.
Thus this step splits after the specified shift, and every resulting
root satisfies $x_{i+1}\equiv\alpha^2/x_{i-1}$ modulo integral series.

For any earlier small variable $x_j$, $j>1$, one has
\begin{equation}\label{eq:pencil-descending-shifts}
 \frac1{x_j}\equiv\frac1{S(x_j)}=S\left(\frac1{x_{j-1}}\right)
 \pmod{k[[t]]}.
\end{equation}
Indeed the difference of the first two expressions is
$x_j^{r-2}/(1+x_j^{r-1})$, which is integral.  Since $f(Y)\equiv Y$ at a
pole, Equations~\ref{eq:pencil-first-shift} and
\ref{eq:pencil-descending-shifts} give successively
\[
 x_{i+s}\equiv\frac{\alpha^2}{x_{i-s}}\pmod{k[[t]]},
 \qquad 1\leq s\leq i-1,
\]
for as long as those variables occur in the tower.  Each step in this
list is a splitting step with the indicated shift.  The last principal
part is $\alpha^2/x_1$, a simple pole.  From then on every extension is
the simple-pole extension of Lemma~\ref{lem:pencil-simple-pole}.
This description includes branches for which the tower stops before the
first nonzero residue, during the descending shifts, or after ramification.

At every step the split factors or the single ramified factor account
for the full relative degree $r$.  Therefore, at each base place,
the scalar extension of the degree-$N$ tower to $k((t))$ is exactly the
product of the branch fields just constructed.  No additional local
factor is omitted.

\subsection{Integral lattices and their determinants}

Fix a base place with parameter $t$, and put $R=k[[t]]$.
The product of its branch integral rings is a free $R$-module of rank
$N$, by the bases constructed above.  Express it in coefficient
coordinates relative to Equation~\ref{eq:pencil-power-basis}, and denote
the resulting lattice in $k((t))^N$ by $\Lambda$.  If $C$ is a basis
matrix of this lattice, define its coefficient determinant order by
$\delta=\operatorname{ord}_t\det C$.

We use two elementary facts about this order.  First, if a square matrix
over $R$ has nonzero determinant, invertible row and column operations
reduce it to a diagonal matrix.  To see this, choose an entry of least
order and move it to the first position.  It divides every entry in $R$,
so row and column subtractions clear its column and row.  Continue in
the remaining square matrix.  The diagonal entries are units times
powers of $t$.  It follows that the dimension over $k$ of the cokernel
equals the determinant order.  Second, for a branch $k((z))$ over
$k((t))$, multiplication by a unit has determinant a unit on $k[[z]]$,
whereas multiplication by $z$ has cokernel $k$ and hence determinant
order one.  Thus
\begin{equation}\label{eq:pencil-norm-order}
 \operatorname{ord}_t\operatorname{Norm}_{k((z))/k((t))}(z^a u)=a
 \qquad(a\in\Z,\ u\in k[[z]]^*).
\end{equation}
The equality for negative $a$ follows by multiplicativity.  The residue
field of each branch is $k$; this is why no residue-degree multiplier
appears in Equation~\ref{eq:pencil-norm-order}.

For completeness, if $b_1,\ldots,b_n$ is a basis of a separable algebra
over a field and $c_1,\ldots,c_r$ is a relative basis for a separable
algebra above it, their composed basis satisfies
\begin{equation}\label{eq:pencil-disc-tower}
 \operatorname{disc}(b_i c_j)
 =\operatorname{disc}(b_i)^r
   \operatorname{Norm}\bigl(\operatorname{disc}(c_j)\bigr).
\end{equation}
Indeed, the evaluation map first applies $r$ copies of the lower
evaluation matrix and then, over each lower embedding, its relative
evaluation matrix.  Taking determinants and squaring gives the formula.
For a product algebra choose the relative bases componentwise; the same
factorization applies.

The discriminant of the power basis in
Equation~\ref{eq:pencil-power-basis} is a nonzero constant.  At every
relative step the roots of the additive equation differ by nonzero
elements of $\F_r$, so its power-basis Vandermonde discriminant is a
nonzero constant; Equation~\ref{eq:pencil-disc-tower} preserves this
property.  For integral bases, a splitting step has relative
discriminant one, using the idempotent basis of its product of rings.
A simple-pole step has relative discriminant order $2(r-1)$ by
Lemma~\ref{lem:pencil-simple-pole}.

Suppose the current local algebra has degree $n$, has $B$ branches,
and its integral discriminant has base order $d_{\rm int}$.  If $R_0$
of those branches ramify at the next step, then
\begin{align*}
 d_{\rm int,new}&=r d_{\rm int}+2(r-1)R_0,\\
 B_{\rm new}&=rB-(r-1)R_0,\qquad n_{\rm new}=rn.
\end{align*}
Equation~\ref{eq:pencil-norm-order} justifies each contribution from a
branch that may already be ramified over the base.  Starting from
$d_{\rm int}=0$ and $n=B=1$, these recurrences give
$d_{\rm int}=2(n-B)$ at every level.  On changing from the power basis
to an integral basis, the trace matrix changes to $C^{\mathsf T}GC$,
where $G$ is the power-basis trace matrix.  Since $\det G$ is a unit,
we obtain
\begin{equation}\label{eq:pencil-lattice-det}
 \delta=N-B,\qquad 0\leq\delta\leq N.
\end{equation}
At a good finite point $a\notin\F_r$ all variables are integral and
all extensions split.  The power basis is then itself an integral
basis: its lattice is contained in the integral lattice, and both have
determinant order zero.

We next give a bound that makes the lattice conditions finite.
At every exceptional place, and in every partial branch, each nonzero
$x_i-c$, $c\in\F_r$, has absolute order at most $N^2$ in the current
parameter.  Before a later ramified extension, the zero-branch orders
and the descending pole orders in the preceding subsection are at
most $N$; the order of $x_i-c$ is zero unless its residue equals $c$
or it has a pole.  Subsequent ramification multiplies these orders
by a total factor at most $N$.  This proves the bound for every such
factor, including variables retained from an earlier level.
Set
\begin{equation}\label{eq:pencil-FH}
 F_0=(l+1)rN^2,\qquad H=(2N-1)F_0+D.
\end{equation}
In every branch, $t^{F_0}e_{\boldsymbol a}$ is integral: the pole order
of the monomial is at most $(l-1)(r-1)N^2$, while $t$ has positive
integral order.  Hence $t^{F_0}R^N\subseteq\Lambda$.

Multiplication by an integral element acts on the integral lattice,
so its trace belongs to $R$.  Each entry of $G$ therefore has order at
least $-2F_0$.  If $y$ is integral, then
$\operatorname{tr}(y e_{\boldsymbol a})$ has order at least $-F_0$;
if $t^D y$ is integral, the lower bound is $-F_0-D$.
The vector of these traces is $G$ times the coefficient vector of $y$.
Since $\det G$ has order zero, each cofactor has order at least
$-2F_0(N-1)$, and Cramer's rule yields
\begin{equation}\label{eq:pencil-lattice-sandwich}
 \begin{aligned}
 t^{F_0}R^N&\subseteq\Lambda_a\subseteq t^{-H}R^N
       &&(a\in\F_r),\\
 t^{F_0-D}R^N&\subseteq t^{-D}\Lambda_\infty
                       \subseteq t^{-H}R^N.
 \end{aligned}
\end{equation}
All shifts here are powers of the base parameter, including the shift
at infinity.  In particular, membership in any target lattice inside
$t^{-H}R^N$ is determined by the coefficient jet with exponents from
$-H$ through $F_0-1$.  Changing the omitted tail adds an element of
the target lattice.  At infinity the smaller upper endpoint
$F_0-D-1$ would also suffice.

\subsection{A large space of global functions}

In coefficient coordinates, let $V_H$ consist of functions whose every
coefficient is a linear combination of
\[
 1,x_1,\ldots,x_1^H,
 \qquad (x_1-a)^{-j}\quad(a\in\F_r, 1\leq j\leq H).
\]
These rational functions are independent: their principal parts at
the distinct finite poles first determine all the negative-power
coefficients, and the remaining polynomial determines the others.
Consequently
\[
 \dim_k V_H=N[1+(r+1)H].
\]
At each of the $r$ bad finite points and at infinity, its coefficient
expansion belongs to $t^{-H}k[[t]]^N$.
Let $\cL_k$ be the subspace of $V_H$ that is integral at every finite
place and becomes integral at infinity after multiplication by $t^D$.

The finite-place quotient by $\Lambda_a$ has dimension
$NH+\delta_a$: scale the inclusion into an integral matrix and apply
the diagonalization argument preceding
Equation~\ref{eq:pencil-norm-order}.  At infinity the target is
$t^{-D}\Lambda_\infty$, whose determinant order is
$\delta_\infty-ND$, because all $N$ coefficient coordinates are scaled.
The quotient dimension there is therefore
$NH+\delta_\infty-ND$.
Expansion followed by these quotient maps gives one linear map from
$V_H$ to their direct sum.  Its rank is at most the dimension of its
target; no independent prescription of local jets is required.
Equation~\ref{eq:pencil-lattice-det} gives
\begin{align}
 \dim_k\cL_k
 &\geq N[1+(r+1)H]
       -\left((r+1)NH+\sum_{a\in\F_r\cup\{\infty\}}\delta_a-ND\right)
       \notag\\
 &\geq N(D-r)=N.\label{eq:pencil-dimension}
\end{align}
There are no conditions to impose elsewhere, since the rational
coefficients and the power basis are integral at every other finite
place.  The space contains the constant function one.

This space is defined over $F$.  Indeed, coefficientwise $q$-Frobenius,
fixing the formal variables $x_i$, preserves the tower equations and
permutes all branches above each exceptional place.  It thus preserves
the subspace $\cL_k$ in the fixed finite rational-coefficient basis of
$V_H$.  Put a matrix whose rows form a basis of this subspace into
reduced row echelon form.  Applying Frobenius gives another reduced row
echelon basis of the same subspace with the same pivots.  Uniqueness of
that form forces every entry to be fixed by Frobenius, hence to belong
to $F$.  Thus the space $\cL$ of its $F$-rational coefficient vectors
has
\begin{equation}\label{eq:pencil-descent}
 \dim_F\cL=\dim_k\cL_k\geq N.
\end{equation}
The algorithm below obtains this space by a finite ground-field
matrix, not by searching for Frobenius-fixed vectors.

\subsection{Evaluation and the quadratic zero bound}

There are $(q-r)N$ rational evaluation tuples above
$x_1\in F\setminus\F_r$.  To verify this count, if
$\beta\in F\setminus\F_r$, then
$\beta^r+\beta\in\F_r^*$ and $\beta^{r+1}\in\F_r^*$, so
$f(\beta)\in\F_r^*$.  The map $S:F\to\F_r$ is linear over $\F_r$
with kernel $\F_r$, hence is onto.  Each equation
$S(Y)=f(\beta)$ has exactly $r$ distinct roots in
$F\setminus\F_r$.  Induction proves the count and describes all tuples.
All functions in $\cL$ are regular at these tuples, so evaluation is
well defined.

Let $a$ be a nonzero element of the span of products of two elements
of $\cL_k$.  It is integral at every finite place, and $t^{2D}a$ is
integral at infinity.  Its field norm to $k(x_1)$ is nonzero.  At every
finite place its local multiplication matrix is integral, so the norm
has no finite poles and is a polynomial in $x_1$.  At infinity,
\[
 0\leq\operatorname{ord}_t\operatorname{Norm}(t^{2D}a)
   =2DN+\operatorname{ord}_t\operatorname{Norm}(a),
\]
and therefore
\begin{equation}\label{eq:pencil-norm-degree}
 \deg\operatorname{Norm}(a)\leq2DN.
\end{equation}
At a good base value $b$, the local algebra is a product of $N$ copies
of $k((x_1-b))$.  The local norm is the product of the integral branch
values $a_1,\ldots,a_N$, and
\[
 \operatorname{ord}_{x_1-b}\operatorname{Norm}(a)
   =\sum_{j=1}^N\operatorname{ord}_{x_1-b}(a_j).
\]
Each zero evaluation contributes at least one to this sum.  In
particular, several zero branches above the same $b$ are counted with
their full multiplicity.  Equation~\ref{eq:pencil-norm-degree} bounds
the total number of zero evaluation positions by $2DN$.

Evaluation on these positions is injective on $\cL$: if a nonzero
function vanished everywhere, its product with the constant one would
contradict $2DN<(q-r)N$, which follows from
Equation~\ref{eq:pencil-parameters}.  Choose $l$ least such that
$N=r^{l-1}\geq d$, and choose $d$ independent functions in $\cL$.
Their $d$ by $(q-r)N$ evaluation matrix has rank $d$.  Choose $d$
independent columns and left-multiply by the inverse of their square
matrix.  The new functions still belong to $\cL$, and their evaluation
vectors include all standard basis vectors.

Call the resulting columns $u_1,\ldots,u_m$.  If a homogeneous
quadratic $Q$ is nonzero on this list, applying $Q$ to the $d$ functions
gives a nonzero element of the product span.  The preceding zero bound
and Equation~\ref{eq:pencil-parameters} give
\[
 \#\{j:Q(u_j)=0\}\leq2DN\leq\lambda_0(q-r)N=\lambda_0m.
\]
Since $N<rd$ for this choice of $l$, we have $m=O_r(d)$.
This proves the algebraic assertions of Lemma~\ref{lem:pencil}; it remains
to justify the claimed algorithm and running time.

\subsection{A finite algorithm and its precision bounds}

The fixed fields $F$ and $\F_r$ can be constructed once and for all.
For example, the roots of $Z^{2^s}-Z$ form the field of $2^s$ elements,
and a representation by a binary basis and an irreducible polynomial
can be fixed by a finite search when $s$ is fixed.  Only the extension
degrees below grow with $d$; they will be constructed by linear algebra.
Field arithmetic in a binary basis has polynomial bit complexity in
the basis dimension.

\paragraph{Finite constant fields.}
At a splitting step, suppose the constant term of the shifted right
side lies in a current finite field $\F_Q$ containing $F$.  If
$S(\beta)$ belongs to that field, then
$\beta^Q-\beta\in\F_r$, because applying $S$ gives zero.  This difference
is fixed by $Q$-Frobenius, so in characteristic two
$\beta^{Q^2}=\beta$.  Thus a quadratic extension suffices to solve the
constant equation.  It can be built without enumerating the field:
the binary-linear map $y\mapsto y^2+y$ on $\F_Q$ has kernel
$\{0,1\}$ and image of codimension one.  Linear algebra finds a vector
$b$ outside that image, and $Y^2+Y+b$ is irreducible and defines the
quadratic extension.  The equation $S(\beta)=b_0$ is itself a binary
linear system in that extension; all its other solutions differ by
elements of $\F_r$.

At most one such quadratic extension is needed at each tower level.
We may work in nested fields, one for each level, and finally express
all constants in a common field of degree
\begin{equation}\label{eq:pencil-constant-degree}
 E=2^{l-1}\leq N
\end{equation}
over $F$.  Simple-pole reversion introduces no new constants.
Thus the use of $k$ in the dimension calculation does not require an
algorithm to enumerate an algebraic closure.

\paragraph{Elementary precision rules.}
Put $M_0=N^2$.  Suppose a nonzero factor $x_i-c$ has order in
$[-M_0,M_0]$ and is known modulo the current parameter to the power $K$,
where $K>M_0$.  Its relative error has order at least $K-M_0$.
Multiplication, integer powers, and inversion preserve this lower
bound on relative accuracy: for inversion use
$(a(1+u))^{-1}=a^{-1}(1+u)^{-1}$ with $u$ of positive order.
A product or quotient of at most a fixed number, depending on $r$,
of these factors therefore has absolute accuracy at least
$K-C_rM_0$, for a fixed constant $C_r$.
The functions $f$, the indicated shifts $\alpha^2/x_j$, and their
images under $S$ are sums of such expressions; one may factor
$Y^{r-1}+1$ over $\F_r^*$ when computing $f$.

Subtraction preserves absolute accuracy even when all the known
principal parts cancel.  In a splitting step, the branch description
already proves that $f(x_i)-S(g)$ is integral.  Computing its input
factors through $L+C_rM_0$ therefore determines it through the
required order $L$.  After the constant equation has been solved,
Equation~\ref{eq:pencil-additive-recursion} determines positive
coefficients through $L$ using only those right-side coefficients.
Adding back $g$ has the same finite accuracy bound.

In a simple-pole step, compute $q_0=1/f(x_i)$, which has order one.
Its inverse under composition through any order uses only its
coefficients through that order and divides by its nonzero linear
coefficient.  Substitute $w^r/(1+w^{r-1})$ in this inverse to obtain
Equation~\ref{eq:pencil-reparameterization}.  If $z(w)$ is known modulo
$w^B$, its relative error has order at least $B-r$.  An old Laurent
series whose pole order in $z$ is at most $M_0$ then changes, after
substitution, only in orders at least $B-r-rM_0$.
Indeed, for each retained term $a_nz^n$ the relative error is at least
$B-r$, and its order is at least $-rM_0$; sums preserve the resulting
absolute bound.  Thus $B\geq L+r(M_0+1)$ suffices for output precision
$L$.  Old coefficients through $\ceil{L/r}+M_0+2$ suffice as well,
since higher powers of $z(w)$ cannot contribute below order $L$.
Updating all previously constructed variables increases the number of
operations, but not the exponent in these precision bounds.

Choose an integer $a_r$ depending only on $r$ so large that
$K_N=(2N)^{a_r}$ dominates $M_0$ and all fixed constants in the preceding
rules, for every $N\geq1$.  The rules are all covered by the single
backward precision budget
\begin{equation}\label{eq:pencil-precision-recurrence}
 L_{i-1}=K_N\bigl(L_i+K_N+H+D\bigr).
\end{equation}
This is an allocated upper budget; the algorithm may retain all
coefficients up to that budget.  It is uniform over levels, branches,
and rational coefficient basis inputs.

\paragraph{The final integrality tests.}
Consider one rational coefficient basis element in $V_H$ times one
power-basis monomial.  In a final branch its pole order, allowing also
the infinity adjustment, is bounded by
\begin{equation}\label{eq:pencil-product-pole}
 P_0=HN+(r-1)(l-1)N^2+DN.
\end{equation}
The first term bounds a base pole of order at most $H$ after total
ramification at most $N$; the second bounds the monomial; the last
is a harmless bound for the infinity multiplier.
Such an input is a product of powers and inverses of the elementary
factors already considered.  Knowing each factor with relative error
of order greater than $P_0$ makes the unknown remainder of the full
input integral.  A sufficiently large fixed multiple
\begin{equation}\label{eq:pencil-final-precision}
 L_{\rm end}=C'_r(H+D+l+1)N^2
\end{equation}
as absolute factor precision therefore determines every negative
coefficient, including after multiplication by $t^D$ at infinity.
Cancellation in a linear combination of these inputs can remove
computed negative coefficients but cannot reveal an uncomputed one:
each individual unknown remainder is already integral.

Use Equation~\ref{eq:pencil-precision-recurrence} backwards from
$L_{\rm end}$, and then perform the forward branch expansions at those
precisions.  The branch type and its prescribed shift are determined
by the explicit residue and descending-shift rules above.  For each
basis input, record all negative coefficients in every branch, with
the infinity multiplier included.  Their vanishing is exactly the
desired system defining $\cL$.

\paragraph{Size and running time.}
There are at most $(r+1)N$ exceptional branches at any level.  Since
$l=1+\log_r N$, Equation~\ref{eq:pencil-FH} gives
$F_0=O_r(N^3)$ and $H=O_r(N^4)$.  The number of rational coefficient
unknowns is $N[1+(r+1)H]=O_r(N^5)$, and
$L_{\rm end}=O_r(N^6)$.  Even retaining a full window of this latter
length per branch gives only $O_r(N^7)$ negative-coefficient equations
over the finite constant fields.  Expanding each coefficient in an
$F$-basis of the common field in
Equation~\ref{eq:pencil-constant-degree} gives at most $O_r(N^8)$
equations over $F$.  Solving this explicit matrix computes $\cL$;
Equation~\ref{eq:pencil-descent} proves that its dimension is at least
$N$.  The computation does not infer dimension by testing random
functions.

Iterating Equation~\ref{eq:pencil-precision-recurrence} for at most
$l-1$ steps shows that every retained series length is at most
\[
 (2N)^{O_r(l+1)}.
\]
Elementary truncated addition, multiplication, inversion and
composition, and coefficientwise reversion, use a number of field
operations polynomial in the retained lengths.  All branch counts,
field representation sizes and matrix dimensions are polynomial in
$N$; allowing every level and every basis input multiplies these
bounds by polynomial factors.  For $N\geq2$,
$\log((2N)^{O_r(l+1)})=O_r((\log N)^2)=O_r(N)$; the case $N=1$ is a
fixed computation.  Total bit complexity is consequently $2^{O_r(N)}$.
Enumerating the rational evaluation tuples, evaluating the functions,
and selecting the independent columns require only additional
polynomial time in $N$ and the matrix sizes.
Finally $N<rd$, so the entire construction takes $2^{O_r(d)}$ time.
Together with the evaluation argument, this completes the proof of
Lemma~\ref{lem:pencil}.
